% GENERATED FILE. Edit canonical lessons, then run npm run generate:books.
% canonical-source-hash: fnv1a64:21577d388aec2519
% canonical-lessons: TA-C402-caman, TA-C402-cutkes, TA-C402-muttai, TA-C402-cuma, TA-C402-irakku

\chapter{Bags and Bundles}
\label{ch:ta-bags-and-bundles}
\hlchaptermodality{402}

\begin{chapteropening}
\textbf{By the end of this chapter:} \emph{I can talk about luggage, a suitcase and a sack, and carrying and unloading them.}

Complete the last lesson of chapter 402: I can talk about luggage, a suitcase and a sack, and carrying and unloading them.
\end{chapteropening}

\section[\texorpdfstring{\ta{சாமான்} (cāmāṉ)}{சாமான் (cāmāṉ)}]{\ta{சாமான்} (cāmāṉ) — luggage; things, goods}

\label{lesson:TA-C402-caman}



\begin{quote}
Before the new one: say the Tamil for to stagger, then the Tamil for a ride.
\end{quote}

\subsection*{You'll want to know: \ta{சாமான்}}
\textbf{\ta{சாமான்}} — \emph{cāmāṉ} — \textquotedblleft{}luggage; things, goods\textquotedblright{}.

\begin{quote}
\ta{என்} \ta{சாமான்கள்} \ta{எல்லாம்} \ta{காரில்} \ta{இருக்கின்றன}.
\end{quote}

\begin{quote}
\emph{eṉ cāmāṉkaḷ ellām kāril irukkiṉṟaṉa.}
\end{quote}

\begin{quote}
All my luggage is in the car.
\end{quote}

Plural: \ta{சாமான்கள்}, \emph{cāmāṉkaḷ}.

Often plural: \ta{சாமான்கள்}, one's things.

\begin{cousinweb}
\textbf{\ta{சாமான்}} (\emph{cāmāṉ}) is a loan from Persian \emph{sāmān}, \textquotedblleft{}goods, equipment\textquotedblright{}, through Hindustani.
\end{cousinweb}

\begin{grammarlens}[title={What you've built}]
One more word for getting around.
\end{grammarlens}

\subsection*{Your turn}
\begin{itemize}
\raggedright
  \item \emph{Say it:} \emph{cāmāṉ}
  \item \emph{Say it:} \emph{cāmāṉ}, once more
  \item \emph{Say it:} \emph{eṉ cāmāṉkaḷ ellām kāril irukkiṉṟaṉa.} — All my luggage is in the car.
  \item \emph{Recall:} say the Tamil for to stagger, then the Tamil for a ride, then say \emph{cāmāṉ} again
\end{itemize}

\begin{tcolorbox}[breakable,colback=teal!4,colframe=teal!35!black,title={Before you move on}]
What does \ta{சாமான்} mean? (\textquotedblleft{}Luggage; things, goods\textquotedblright{}.) Say it once more.
\end{tcolorbox}
\section[\texorpdfstring{\ta{சூட்கேஸ்} (cūṭkēs)}{சூட்கேஸ் (cūṭkēs)}]{\ta{சூட்கேஸ்} (cūṭkēs) — a suitcase}

\label{lesson:TA-C402-cutkes}



\begin{quote}
Before the new one: say the Tamil for a ride, then the Tamil for luggage.
\end{quote}

\subsection*{You'll want to know: \ta{சூட்கேஸ்}}
\textbf{\ta{சூட்கேஸ்}} — \emph{cūṭkēs} — \textquotedblleft{}a suitcase\textquotedblright{}.

\begin{quote}
\ta{அவள்} \ta{சூட்கேஸை} \ta{மெதுவாகத்} \ta{தூக்கினாள்}.
\end{quote}

\begin{quote}
\emph{avaḷ cūṭkēsai metuvākat tūkkiṉāḷ.}
\end{quote}

\begin{quote}
She lifted the suitcase slowly.
\end{quote}

A loan in native letters: before -\ta{ஐ} the final \ta{ஸ்} takes the vowel, \ta{சூட்கேஸை}.

\begin{cousinweb}
\textbf{\ta{சூட்கேஸ்}} (\emph{cūṭkēs}) is a loan from English \emph{suitcase}, written with the grantha letter \ta{ஸ}.
\end{cousinweb}

\begin{grammarlens}[title={What you've built}]
One more word for getting around.
\end{grammarlens}

\subsection*{Your turn}
\begin{itemize}
\raggedright
  \item \emph{Say it:} \emph{cūṭkēs}
  \item \emph{Say it:} \emph{cūṭkēs}, once more
  \item \emph{Say it:} \emph{avaḷ cūṭkēsai metuvākat tūkkiṉāḷ.} — She lifted the suitcase slowly.
  \item \emph{Recall:} say the Tamil for a ride, then the Tamil for luggage, then say \emph{cūṭkēs} again
\end{itemize}

\begin{tcolorbox}[breakable,colback=teal!4,colframe=teal!35!black,title={Before you move on}]
What does \ta{சூட்கேஸ்} mean? (\textquotedblleft{}A suitcase\textquotedblright{}.) Say it once more.
\end{tcolorbox}
\section[\texorpdfstring{\ta{மூட்டை} (mūṭṭai)}{மூட்டை (mūṭṭai)}]{\ta{மூட்டை} (mūṭṭai) — a sack; a bundle}

\label{lesson:TA-C402-muttai}



\begin{quote}
Before the new one: say the Tamil for luggage, then the Tamil for a suitcase.
\end{quote}

\subsection*{You'll want to know: \ta{மூட்டை}}
\textbf{\ta{மூட்டை}} — \emph{mūṭṭai} — \textquotedblleft{}a sack; a bundle\textquotedblright{}.

\begin{quote}
\ta{கடையின்} \ta{முன்னால்} \ta{அரிசி} \ta{மூட்டைகள்} \ta{இருந்தன}.
\end{quote}

\begin{quote}
\emph{kaṭaiyiṉ muṉṉāl arici mūṭṭaikaḷ iruntaṉa.}
\end{quote}

\begin{quote}
There were sacks of rice in front of the shop.
\end{quote}

Also a measure: \ta{ஒரு} \ta{மூட்டை} \ta{அரிசி}, a sack of rice.

\begin{cousinweb}
\textbf{\ta{மூட்டை}} (\emph{mūṭṭai}) is native Dravidian (DEDR 5037), \textquotedblleft{}a bundle, something tied up, a load carried\textquotedblright{}.
\end{cousinweb}

\begin{grammarlens}[title={What you've built}]
One more word for getting around.
\end{grammarlens}

\subsection*{Your turn}
\begin{itemize}
\raggedright
  \item \emph{Say it:} \emph{mūṭṭai}
  \item \emph{Say it:} \emph{mūṭṭai}, once more
  \item \emph{Say it:} \emph{kaṭaiyiṉ muṉṉāl arici mūṭṭaikaḷ iruntaṉa.} — There were sacks of rice in front of the shop.
  \item \emph{Recall:} say the Tamil for luggage, then the Tamil for a suitcase, then say \emph{mūṭṭai} again
\end{itemize}

\begin{tcolorbox}[breakable,colback=teal!4,colframe=teal!35!black,title={Before you move on}]
What does \ta{மூட்டை} mean? (\textquotedblleft{}A sack; a bundle\textquotedblright{}.) Say it once more.
\end{tcolorbox}
\section[\texorpdfstring{\ta{சும} (cuma)}{சும (cuma)}]{\ta{சும} (cuma) — to carry; to bear a load}

\label{lesson:TA-C402-cuma}



\begin{quote}
Before the new one: say the Tamil for a suitcase, then the Tamil for a sack.
\end{quote}

\subsection*{You'll want to know: \ta{சும}}
\textbf{\ta{சும}} — \emph{cuma} — \textquotedblleft{}to carry; to bear a load\textquotedblright{}.

\begin{quote}
\ta{நான்} \ta{அந்த} \ta{மூட்டையை} \ta{வீடு} \ta{வரை} \ta{சுமந்தேன்}.
\end{quote}

\begin{quote}
\emph{nāṉ anta mūṭṭaiyai vīṭu varai cumantēṉ.}
\end{quote}

\begin{quote}
I carried that sack all the way home.
\end{quote}

Past (I …): \ta{சுமந்தேன்}, \emph{cumantēṉ}.

The noun is \ta{சுமை}, a load; \ta{சும} is its verb, with the past in -\ta{ந்த}-: \ta{சுமந்தேன்}.

\begin{cousinweb}
\textbf{\ta{சும}} (\emph{cuma}) is native Dravidian (DEDR 2677); the same article holds \ta{சுமை}, \textquotedblleft{}a burden, a load\textquotedblright{}.
\end{cousinweb}

\begin{grammarlens}[title={What you've built}]
One more word for getting around.
\end{grammarlens}

\subsection*{Your turn}
\begin{itemize}
\raggedright
  \item \emph{Say it:} \emph{cuma}
  \item \emph{Say it:} \emph{cuma}, once more
  \item \emph{Say it:} \emph{nāṉ anta mūṭṭaiyai vīṭu varai cumantēṉ.} — I carried that sack all the way home.
  \item \emph{Recall:} say the Tamil for a suitcase, then the Tamil for a sack, then say \emph{cuma} again
\end{itemize}

\begin{tcolorbox}[breakable,colback=teal!4,colframe=teal!35!black,title={Before you move on}]
What does \ta{சும} mean? (\textquotedblleft{}To carry; to bear a load\textquotedblright{}.) Say it once more.
\end{tcolorbox}
\section[\texorpdfstring{\ta{இறக்கு} (iṟakku)}{இறக்கு (iṟakku)}]{\ta{இறக்கு} (iṟakku) — to unload; to take down}

\label{lesson:TA-C402-irakku}



\begin{quote}
Before the new one: say the Tamil for a sack, then the Tamil for to carry.
\end{quote}

\subsection*{You'll want to know: \ta{இறக்கு}}
\textbf{\ta{இறக்கு}} — \emph{iṟakku} — \textquotedblleft{}to unload; to take down\textquotedblright{}.

\begin{quote}
\ta{ஓட்டுநர்} \ta{சாமான்களை} \ta{வண்டியிலிருந்து} \ta{இறக்கினார்}.
\end{quote}

\begin{quote}
\emph{ōṭṭunar cāmāṉkaḷai vaṇṭiyiliruntu iṟakkiṉār.}
\end{quote}

\begin{quote}
The driver unloaded the luggage from the vehicle.
\end{quote}

Past (I …): \ta{இறக்கினேன்}, \emph{iṟakkiṉēṉ}.

\ta{இறங்கு} is to get down oneself; \ta{இறக்கு} is to bring something down.

\begin{cousinweb}
\textbf{\ta{இறக்கு}} (\emph{iṟakku}) is native Dravidian (DEDR 516), in the same article as \ta{இறங்கு}, \textquotedblleft{}to descend\textquotedblright{}.
\end{cousinweb}

\begin{grammarlens}[title={What you've built}]
One more word for getting around.
\end{grammarlens}

\subsection*{Your turn}
\begin{itemize}
\raggedright
  \item \emph{Say it:} \emph{iṟakku}
  \item \emph{Say it:} \emph{iṟakku}, once more
  \item \emph{Say it:} \emph{ōṭṭunar cāmāṉkaḷai vaṇṭiyiliruntu iṟakkiṉār.} — The driver unloaded the luggage from the vehicle.
  \item \emph{Recall:} say the Tamil for a sack, then the Tamil for to carry, then say \emph{iṟakku} again
\end{itemize}

\begin{tcolorbox}[breakable,colback=teal!4,colframe=teal!35!black,title={Before you move on}]
What does \ta{இறக்கு} mean? (\textquotedblleft{}To unload; to take down\textquotedblright{}.) Say it once more.
\end{tcolorbox}
